\reading{CR-9}{Estimating Default Probabilities}{STRIKE FIRST}{8}

\begin{crkeybox}[Learning objectives for this reading]
\begin{enumerate}[label=\textbf{\alph*.},leftmargin=2em]
\item Compare agencies' ratings to internal credit rating systems.
\item Describe Altman's Z-score methodology and apply it to predict a firm's default and to classify a sample of firms by credit quality.
\item Describe the relationship between borrower rating and probability of default.
\item Describe a rating migration matrix and calculate the probability of default, cumulative probability of default, and marginal probability of default.
\item Define the hazard rate and use it to define probability functions for default time as well as to calculate conditional and unconditional default probabilities.
\item Describe recovery rates and their dependencies on default rates.
\item Define a credit default swap (CDS) and explain its mechanics including the obligations of both the default protection buyer and the default protection seller.
\item Describe CDS spreads and explain how CDS spreads can be used to estimate hazard rates.
\item Define and explain CDS-bond basis.
\item Compare default probabilities calculated from historical data with those calculated from credit yield spreads.
\item Describe the difference between real-world and risk-neutral default probabilities and determine which one to use in the analysis of credit risk.
\item Calculate the value of a firm's debt and equity, the volatility of firm value, and the volatility of firm equity using the Merton model.
\item Calculate distance to default and default probability using the Merton model.
\item Assess the quality of the default probabilities produced by the Merton model, the Moody's KMV model, and the Kamakura model.
\end{enumerate}
\end{crkeybox}

% =====================================================================
\lo{CR-9 a}{Compare agencies' ratings to internal credit rating systems.}

Credit ratings provided by rating agencies assess the credit quality of larger
publicly traded bond issuers, indicating their ability to repay debt.

\begin{center}
\footnotesize
\begin{tabular}{L{6.9cm} L{7.6cm}}
\toprule
\thead{Agencies' ratings (large firms)} & \thead{Internal credit rating systems (SMEs)} \\
\midrule
Provided by established agencies: Moody's, S\&P, Fitch &
  Major agencies typically do not cover small and mid-sized firms, so lenders
  rely on internally developed systems \\
Publicly available information used for assessment &
  Built by the lender using financial statement metrics \\
\textbf{Stability in ratings}; agencies adopt a long-term view on
creditworthiness &
  Used for estimating probabilities of default directly \\
\textbf{Avoiding reversals} &
  Factors: profitability (e.g., return on equity and return on assets); liquidity
  (e.g., quick ratio); solvency (e.g., debt-to-equity); balance sheet ratios \\
Focus on \textbf{long-term changes in creditworthiness}: issuer performance
through economic cycles, to mitigate short-term volatility's impact on ratings &
  \textbf{Importance of cash in loan repayment}: since repayment relies on cash,
  lenders analyse the cash flow statement to assess repayment ability \\
\bottomrule
\end{tabular}
\end{center}
\figcap{The contrast that carries the weight is deliberate stability versus
direct estimation. Agencies rate through the cycle and avoid reversals; an
internal system is built to produce a PD number the lender can use now.}

% =====================================================================
\lo{CR-9 b}{Describe Altman's Z-score methodology and apply it to predict a firm's default and to classify a sample of firms by credit quality.}

\subsection*{Linear discriminant analysis}

\term{Linear discriminant analysis (LDA)} is a statistical method used for
classification and dimensionality reduction. Its main goal is to find a linear
combination of features that characterizes or separates two or more classes in
the data.

A \term{scoring function} is a linear function of variables produced by an LDA.
The variables are chosen based on their estimated contribution to the likelihood
of default, and come from an extensive pool of qualitative features and
accounting ratios.

LDA categorizes companies into two groups: those that are doing well
(\term{solvent}) and those that are struggling (\term{insolvent}). The challenge
is to predict which companies will become insolvent \emph{before} it happens. To
do this, the model assigns an Altman Z-score to each company.

\subsection*{Altman's Z-score}

The Z-score predicts the likelihood of a company going into bankruptcy from five
accounting ratios.

\begin{crfmlbox}[Altman's Z-score, original model for publicly traded manufacturing companies]
\[ Z \;=\; 1.2X_1 + 1.4X_2 + 3.3X_3 + 0.6X_4 + 0.999X_5 \]
\vk{$X_1$ = working capital / total assets; $X_2$ = retained earnings / total
assets; $X_3$ = earnings before interest and taxes / total assets; $X_4$ = market
value of equity / book value of total liabilities; $X_5$ = sales / total assets.
Note that $X_4$ is the only ratio whose denominator is not total assets, and the
only one that uses a \emph{market} value in the numerator.}
\end{crfmlbox}

The higher the Z-score, the better (more solvent). The Z-score cut-off, also
known as the \term{discriminant threshold}, was set at $Z = 2.675$.

\begin{center}
\footnotesize
\begin{tabular}{L{4.0cm} L{7.0cm}}
\toprule
\thead{Z-score} & \thead{Interpretation} \\
\midrule
Greater than 3.0 & Unlikely to default \\
Between 2.7 and 3.0 & On alert \\
Between 1.8 and 2.7 & Good chance of default \\
Less than 1.8 & Very high probability of financial embarrassment \\
\bottomrule
\end{tabular}
\end{center}

\begin{crexbox}[Altman Z-score]
Consider a company for which working capital is 170,000, total assets are
670,000, earnings before interest and taxes is 60,000, sales are 2,200,000, the
market value of equity is 380,000, total liabilities is 240,000, and retained
earnings is 300,000.
\tcblower
\begin{align*}
X_1 &= 170{,}000 / 670{,}000 = 0.254 \\
X_2 &= 300{,}000 / 670{,}000 = 0.448 \\
X_3 &= 60{,}000 / 670{,}000 = 0.0896 \\
X_4 &= 380{,}000 / 240{,}000 = 1.583 \\
X_5 &= 2{,}200{,}000 / 670{,}000 = 3.284
\end{align*}
The Z-score is
\[ 1.2 \times 0.254 + 1.4 \times 0.448 + 3.3 \times 0.0896 + 0.6 \times 1.583
   + 0.999 \times 3.284 = 5.46 \]
The Z-score indicates that the company is not in danger of defaulting in the near
future. Note that $X_5$ alone contributes 3.28 of the 5.46: the asset turnover
term dominates the total for a company with high sales relative to its asset
base.
\end{crexbox}

% =====================================================================
\lo{CR-9 c}{Describe the relationship between borrower rating and probability of default.}

\begin{enumerate}
\item \term{Inverse relationship.} Higher borrower ratings indicate a lower
probability of default (AAA), while lower ratings (B) suggest a higher risk of
default.
\item \term{Investment-grade dynamics.} For investment-grade bonds, the
probability of default generally \emph{increases} over the initial years,
reflecting potential financial deterioration over time.
\item \term{Non-investment-grade trends.} Non-investment-grade bonds show an
initial increase in default risk, which may \emph{decrease} over time if the
issuer survives the early critical period. This indicates possible financial
stabilization and an improvement in perceived creditworthiness after surviving
the initial high-risk phase.
\end{enumerate}

% =====================================================================
\lo{CR-9 d}{Describe a rating migration matrix and calculate the probability of default, cumulative probability of default, and marginal probability of default.}

S\&P borrower ratings range from AAA (highest rating, best quality) to D (lowest
rating, in payment default). On this scale, investment-grade bonds are rated from
AAA to BBB and non-investment-grade (junk) bonds from BB to D.

\begin{center}
\footnotesize
\begin{tabular}{l r r r r r r r r}
\toprule
\thead{Rating} & \thead{1} & \thead{2} & \thead{3} & \thead{4} & \thead{5} &
\thead{7} & \thead{10} & \thead{15} \\
\midrule
AAA & 0.00 & 0.03 & 0.13 & 0.24 & 0.34 & 0.51 & 0.70 & 0.90 \\
AA & 0.02 & 0.06 & 0.11 & 0.21 & 0.30 & 0.49 & 0.70 & 0.99 \\
A & 0.05 & 0.13 & 0.22 & 0.33 & 0.46 & 0.76 & 1.20 & 1.84 \\
BBB & 0.16 & 0.43 & 0.75 & 1.14 & 1.54 & 2.27 & 3.24 & 4.54 \\
BB & 0.63 & 1.93 & 3.46 & 4.99 & 6.43 & 8.89 & 11.64 & 14.65 \\
B & 3.34 & 7.80 & 11.75 & 14.89 & 17.35 & 20.99 & 24.62 & 28.24 \\
CCC/C & 28.30 & 38.33 & 43.42 & 46.36 & 48.58 & 50.75 & 52.76 & 54.76 \\
\bottomrule
\end{tabular}
\end{center}
\figcap{Cumulative issuer-weighted default rates for 1981--2020, in \%, by time
horizon in years. Source: S\&P Global Ratings Research. Every figure in this
table is \emph{cumulative}; the marginal and conditional probabilities below are
all derived from differences within a row.}

\begin{crfmlbox}[The three probabilities from one table]
\[ \text{Cumulative PD}(t) \;=\; \text{read directly from the table} \]
\[ \text{Marginal (unconditional) PD}(t) \;=\; \text{Cumulative PD}(t) -
   \text{Cumulative PD}(t-1) \]
\[ \text{Conditional PD}(t) \;=\; \frac{\text{Marginal PD}(t)}
   {\text{Survival probability}(t-1)} \]
\vk{Survival probability$(t-1) = 1 - \text{Cumulative PD}(t-1)$. The marginal
probability is unconditional: it is the chance of defaulting in year $t$ viewed
from today. The conditional probability is the chance of defaulting in year $t$
\emph{given} survival to the start of it, so it is always the larger of the two.}
\end{crfmlbox}

\begin{crexbox}[Marginal probability of default from the matrix]
For a bond with an initial rating of AA, calculate the marginal probability of
default for the third year.
\tcblower
Use the cumulative probabilities for year 2 and year 3:
\[ \mathrm{MPD}(3) = P(3) - P(2) = 0.11\% - 0.06\% = 0.05\% \]
\end{crexbox}

\begin{center}
\begin{tikzpicture}
\pgfplotsset{
  mpd/.style={
    ybar, bar width=9pt, width=7.4cm, height=5.4cm,
    xlabel={Year}, symbolic x coords={1,2,3,4}, xtick=data, ymin=0,
    xlabel style={font=\sffamily\scriptsize}, ylabel style={font=\sffamily\scriptsize},
    tick label style={font=\sffamily\scriptsize},
    axis y line*=left, axis x line*=bottom,
    nodes near coords,
    every node near coord/.append style={font=\sffamily\tiny,
      /pgf/number format/.cd, fixed, precision=2}}}
\begin{axis}[mpd, name=left, ylabel={Marginal PD (\%)}, ymax=5.6,
  title={\sffamily\scriptsize\bfseries\color{FRMRed}B rated}]
\addplot[draw=FRMRed, fill=FRMRed!45, area legend]
  coordinates {(1,3.34) (2,4.46) (3,3.95) (4,3.14)};
\end{axis}
\begin{axis}[mpd, at={(left.south east)}, xshift=2.0cm, anchor=south west,
  ylabel={Marginal PD (\%)}, ymax=0.13, ytick={0,0.05,0.10}, yticklabels={0,0.05,0.10},
  title={\sffamily\scriptsize\bfseries\color{FRMNavy}AA rated}]
\addplot[draw=FRMNavy, fill=FRMNavy!45, area legend]
  coordinates {(1,0.02) (2,0.04) (3,0.05) (4,0.10)};
\end{axis}
\end{tikzpicture}
\end{center}
\figcap{Marginal default probabilities derived from the table above, and the
picture behind objective c. \textbf{The two panels are on different vertical
scales} --- note the axis maxima, which differ by a factor of more than forty ---
because on a shared scale the AA bars are invisible. What the panels are for is
the \emph{shape}. The investment-grade profile climbs steadily. The
non-investment-grade profile peaks in year 2 and then falls, because a B-rated
issuer that survives the early critical period has demonstrably stabilised.}

% =====================================================================
\lo{CR-9 e}{Define the hazard rate and use it to define probability functions for default time as well as to calculate conditional and unconditional default probabilities.}

\begin{crdefbox}[Hazard rate]
The hazard rate, also known as \term{default intensity}, is represented by the
(constant) parameter $\lambda$. The probability of default over the next small
time interval $dt$ is $\lambda\,dt$.
\end{crdefbox}

\begin{crfmlbox}[Probability of default by time $t$]
\[ Q(t) \;=\; 1 - e^{-\bar{\lambda}(t)\times t} \]
\vk{$\bar{\lambda}(t)$ = the \emph{average} hazard rate between time 0 and time
$t$. $Q(t)$ is a cumulative probability. Survival to $t$ is
$e^{-\bar{\lambda}(t)t}$.}
\end{crfmlbox}

\begin{crexbox}[Hazard rate of 1.5\% per year]
Suppose the hazard rate is a constant 1.5\% per year.
\tcblower
Cumulative PD by the end of the first year:
\[ 1 - e^{-0.015 \times 1} = 0.0149 \]
Cumulative PD by the end of the second year:
\[ 1 - e^{-0.015 \times 2} = 0.0296 \]
The cumulative PDs by the end of the third, fourth, and fifth years are similarly
0.0440, 0.0582, and 0.0723.

\smallskip
The \textbf{unconditional} PD during the fourth year:
\[ 0.0582 - 0.0440 = 0.0142 \]
The \textbf{conditional} PD in the fourth year, conditional on no earlier default:
\[ 0.0142 \,/\, (1 - 0.0440) = 0.0149 \]
\end{crexbox}

\begin{crexbox}[Hazard rate of 2\% per year]
Assume a 2\% hazard rate.
\tcblower
The probability of default by the end of Year 4 is $7.69\% = 1 - e^{-0.02\times4}$.

By the end of Year 5 it is $9.52\% = 1 - e^{-0.02\times5}$.

The \textbf{unconditional} probability of default during Year 5 is
$9.52\% - 7.69\% = 1.83\%$.

The \textbf{conditional} probability of default in Year 5, conditional on no
earlier default, is $1.83\% / (1 - 7.69\%) = 1.98\%$.
\end{crexbox}

\begin{center}
\begin{tikzpicture}
\begin{axis}[
  width=12.6cm, height=5.8cm,
  xlabel={Year}, ylabel={Probability of default in that year (\%)},
  xlabel style={font=\sffamily\scriptsize}, ylabel style={font=\sffamily\scriptsize},
  tick label style={font=\sffamily\scriptsize},
  xmin=0.5, xmax=20.5, ymin=1.1, ymax=2.15,
  xtick={1,5,10,15,20},
  grid=major, grid style={FRMRule!50, line width=0.3pt},
  legend style={font=\scriptsize\sffamily, draw=FRMRule, fill=white,
    at={(0.5,-0.32)}, anchor=north, legend columns=-1, column sep=10pt}]
\addplot[domain=1:20, samples=20, draw=FRMNavy, line width=1.1pt, sharp plot,
  mark=*, mark size=1.2pt]
  {100*exp(-0.02*(x-1))*(1-exp(-0.02))};
\addlegendentry{unconditional (marginal) PD in year $t$}
\addplot[domain=1:20, samples=2, draw=FRMGold, line width=1.1pt, sharp plot]
  {100*(1-exp(-0.02))};
\addlegendentry{conditional PD in year $t$}
\addplot[only marks, mark=*, mark size=2pt, draw=FRMRed, fill=FRMRed]
  coordinates {(5,1.8281)};
\node[font=\sffamily\scriptsize\bfseries, text=FRMRed, fill=white, inner sep=1.5pt]
  at (axis cs:8.4,1.60) {$1.83\%$ at $t=5$};
\draw[-{Stealth[length=4pt]}, FRMRed] (axis cs:6.9,1.65) -- (axis cs:5.25,1.80);
\node[font=\sffamily\scriptsize\bfseries, text=FRMGold, fill=white, inner sep=1.5pt]
  at (axis cs:15.5,2.06) {constant at $1.98\%$};
\end{axis}
\end{tikzpicture}
\end{center}
\figcap{The two probabilities under a constant 2\% hazard rate, plotted year by
year. The conditional probability never moves: given survival to the start of any
year, the chance of defaulting in it is always $1-e^{-0.02}=1.98\%$. The
unconditional probability decays, because it is the conditional probability
multiplied by the shrinking chance of having survived that long. They coincide
only in year 1, and the gap widens forever after. The 1.83\% computed in the
example above is the year-5 point on the navy line.}

\subsection*{Approximating hazard rates from credit spreads}

\begin{crfmlbox}[Hazard rate from a credit spread]
\[ \bar{\lambda} \;=\; \frac{s(T)}{1-\mathrm{RR}} \]
\vk{$\bar{\lambda}$ = average hazard rate; $s(T)$ = credit spread for maturity
$T$ (a CDS spread or a bond yield spread); RR = recovery rate. The denominator is
the loss given default, so the spread is being grossed up by the fraction
actually lost on default.}
\end{crfmlbox}

\begin{crexbox}[Hazard rate from a single spread]
Five-year credit spread: 210 basis points (2.1\%). Expected recovery rate: 30\%.
What is the hazard rate?
\tcblower
\[ \bar{\lambda} = \frac{0.021}{1 - 0.30} = \frac{0.021}{0.70} = 0.030 = 3.00\%
   \text{, or 300 basis points} \]
\end{crexbox}

\begin{crexbox}[Bootstrapping hazard rates across maturities]
Credit spreads for 3-year, 5-year, and 10-year CDS are 80 bps, 90 bps, and 110
bps respectively. The expected recovery rate is 65\%. Calculate the average
hazard rate for each maturity, and then the forward hazard rate between year 3
and year 5.
\tcblower
With $1-\mathrm{RR} = 0.35$:
\begin{align*}
\bar{\lambda}(3) &= 0.0080 / 0.35 = 0.0229 \text{, or } 2.29\% \\
\bar{\lambda}(5) &= 0.0090 / 0.35 = 0.0257 \text{, or } 2.57\% \\
\bar{\lambda}(10) &= 0.0110 / 0.35 = 0.0314 \text{, or } 3.14\%
\end{align*}

To find the average hazard rate \emph{between} Year 3 and Year 5, note that the
five-year average must be the time-weighted blend of the first three years and
the next two:
\[ \bar{\lambda}(5)\times 5 \;=\; \bar{\lambda}(3)\times 3
   \;+\; \lambda_{3\to5}\times 2 \]
\[ \lambda_{3\to5} = \frac{(5\times 0.0257) - (3\times 0.0229)}{2} = 0.0299
   \text{, or } 2.99\% \]
Carrying the unrounded averages through gives exactly 3.00\%. The same method
gives the year 5 to year 10 forward rate as
$[(10\times0.0314)-(5\times0.0257)]/5 = 3.71\%$.
\end{crexbox}

% =====================================================================
\lo{CR-9 f}{Describe recovery rates and their dependencies on default rates.}

\begin{crdefbox}[Recovery rate]
The recovery rate (RR) for a bond is the price at which it trades about one month
after default, as a percentage of its principal amount.
\end{crdefbox}

When an issuer files for bankruptcy, most creditors are unlikely to receive full
repayment of the face value of their claims, which is what makes the recovery
rate relevant.

Data from 1983 to 2021 shows senior bonds like first lien bonds have higher
average recovery rates (54.6\%) than junior ones like subordinated bonds
(32.1\%).

\begin{crkeybox}[The direction, and the mechanism behind it]
Recovery rates are \term{inversely} related to default rates. Higher defaults lead
to lower recoveries.

\smallskip
The mechanism, from the 2007--2008 crisis: increasing mortgage default rates lead
to more foreclosures, a surplus of houses for sale, and a decline in house
prices. Consequently recovery rates decrease.

\smallskip
In a robust economy defaults are rarer and recovery rates are higher; in a weak
economy defaults increase and recoveries drop. This dynamic poses a significant
challenge to lenders during downturns, who face high default rates and low
recovery rates \emph{simultaneously}.
\end{crkeybox}

% =====================================================================
\lo{CR-9 g}{Define a credit default swap (CDS) and explain its mechanics including the obligations of both the default protection buyer and the default protection seller.}

\begin{crdefbox}[Credit default swap]
A CDS is a type of credit derivative where one party (the \term{protection
buyer}) pays another party (the \term{protection seller}) regular payments in
exchange for credit protection on a reference entity, such as a corporation. The
essence of a CDS is to transfer the credit exposure of fixed income products
between parties.
\end{crdefbox}

\begin{center}
\begin{tikzpicture}[font=\sffamily\footnotesize]
\node[draw=FRMNavy, fill=PaleNavy, line width=0.8pt, rounded corners=1pt,
      text width=3.7cm, align=center, minimum height=1.9cm] (b) at (0,0)
      {\textbf{Protection buyer}\\[2pt]\scriptsize e.g.\ FI Advisors\\[1pt]\scriptsize holds exposure to\\\scriptsize the reference entity};
\node[draw=FRMGold, fill=PaleGold, line width=0.8pt, rounded corners=1pt,
      text width=3.7cm, align=center, minimum height=1.9cm] (s) at (10.2,0)
      {\textbf{Protection seller}\\[2pt]\scriptsize e.g.\ Market Makers, Inc.\\[1pt]\scriptsize takes on the\\\scriptsize credit exposure};
% premium leg
\draw[-{Stealth[length=6pt]}, FRMNavy, line width=1.1pt]
  ($(b.east)+(0,0.55)$) -- ($(s.west)+(0,0.55)$);
\node[font=\scriptsize\bfseries, text=FRMNavy] at (5.1,1.05) {PREMIUM LEG};
\node[font=\scriptsize, text=FRMGrey] at (5.1,0.2)
  {regular payments of the CDS spread};
% protection leg
\draw[-{Stealth[length=6pt]}, FRMGold, line width=1.1pt]
  ($(s.west)-(0,0.85)$) -- ($(b.east)-(0,0.85)$);
\node[font=\scriptsize\bfseries, text=FRMGold] at (5.1,-0.5) {PROTECTION LEG};
\node[font=\scriptsize, text=FRMGrey] at (5.1,-1.25)
  {pays \emph{only if} a credit event occurs};
% reference entity
\node[draw=FRMRule, fill=white, line width=0.7pt, rounded corners=1pt,
      text width=3.3cm, align=center] (r) at (5.1,-2.85)
      {\scriptsize\textbf{Reference entity}\\\scriptsize e.g.\ ELF};
\draw[FRMGrey, dashed, line width=0.7pt] (b.south) |- (r.west);
\draw[FRMGrey, dashed, line width=0.7pt] (s.south) |- (r.east);
\end{tikzpicture}
\end{center}
\figcap{The two legs of a CDS. The premium leg runs continuously from buyer to
seller; the protection leg is contingent and runs the other way. Neither party
need own the reference entity's debt, and the reference entity is not a party to
the contract.}

\subsection*{Settlement}

\begin{enumerate}
\item \term{Physical delivery.} The buyer delivers the defaulted bonds to the
seller in exchange for the face value.
\item \term{Cash settlement.} The seller pays the buyer the fair market value of
the defaulted bonds, based on the cheapest-to-deliver bond.
\end{enumerate}

\begin{crexbox}[A CDS in use]
FI Advisors (FIA) hedges against ELF's potential default by entering a CDS with
Market Makers (MM), paying \$1.5 million annually. The CDS spread is 75 basis
points.
\tcblower
If ELF does not default, FIA gets no compensation: the premium is simply the cost
of the hedge. In the case of default, MM pays FIA \$200 million, or a loss-based
cash settlement.
\end{crexbox}

\subsection*{Credit indices}

Credit indices are baskets of CDS for multiple companies. They allow investors to
buy or sell protection against defaults of a \emph{group} of issuers, instead of
focusing on individual companies. These indices typically offer CDS contracts
with maturities of 3, 5, 7, and 10 years, with settlement dates likely on June 20
and December 20.

The two well-known investment-grade credit indices each contain \textbf{125}
companies (reference entities): \term{CDX NA IG} (North America) and
\term{iTraxx Europe}.

\begin{crexbox}[Pricing an index position]
The 5-year iTraxx Europe index is quoted as bid 150 basis points and ask 152 bps.
An investor seeks \euro500,000 of protection per company in the index. What are the
annual payments?
\tcblower
Using the ask price, since the investor is buying protection:
\[ 0.0152 \times \text{\euro}500{,}000 \times 125 = \text{\euro}950{,}000 \]
If one reference entity defaults, the investor receives the appropriate payoff as
well as a future reduction in the annual payments of
\[ \text{\euro}950{,}000 \,/\, 125 = \text{\euro}7{,}600 \]
\end{crexbox}

\subsection*{Fixed coupons}

\begin{itemize}
\item \term{Fixed coupon.} A standardized fee of 100 basis points (1\%) per year,
typically paid by the protection buyer for investment-grade issuers. Having a
fixed coupon simplifies the process of clearing and settling CDS contracts,
making them more efficient.
\item \term{CDS spread.} The actual cost of protection, reflecting the perceived
risk of default for the reference entity. It can be higher or lower than the fixed
coupon.
\end{itemize}

The fixed coupon might not match the actual CDS spread, so an up-front premium
reconciles the two:

\begin{center}
\footnotesize
\begin{tabular}{L{5.2cm} L{9.3cm}}
\toprule
\thead{Relationship} & \thead{Up-front premium} \\
\midrule
Spread $=$ fixed coupon & No up-front premium is required \\
Spread $>$ fixed coupon & The \textbf{buyer} pays an up-front premium reflecting
the difference \\
Spread $<$ fixed coupon & The \textbf{seller} pays an up-front premium to the
buyer for the difference. This is less common but can occur \\
\bottomrule
\end{tabular}
\end{center}
\figcap{Who pays the up-front premium is a role-and-direction fact, and both
halves are swappable. The party that is getting the better of the standardised
coupon compensates the other.}

% =====================================================================
\lo{CR-9 h}{Describe CDS spreads and explain how CDS spreads can be used to estimate hazard rates.}

A \term{CDS spread} is essentially the cost of protection against a default on a
specific debt. It is expressed in basis points, as a percentage of the notional
amount of the debt being insured.

The route from a CDS spread to a hazard rate is the approximation given at
objective e: $\bar{\lambda} = s(T) / (1-\mathrm{RR})$. A CDS spread is a credit
spread, so it can be substituted for $s(T)$ directly, and bootstrapped across
maturities as shown there.

% =====================================================================
\lo{CR-9 i}{Define and explain CDS-bond basis.}

\begin{crdefbox}[Bond yield spread and the CDS-bond basis]
The \term{bond yield spread} is the difference between the yield of a corporate
bond and a comparable risk-free bond, effectively measuring the extra yield an
investor earns for taking on additional risk.

\smallskip
The \term{CDS-bond basis} should be zero, or near zero, because the CDS spread
and the bond yield spread should be equal.
\end{crdefbox}

\begin{crfmlbox}[CDS-bond basis]
\[ \text{CDS-bond basis} \;=\; \text{CDS spread} \;-\; \text{bond yield spread} \]
\end{crfmlbox}

\subsection*{Strategy based on spread comparison}

\begin{enumerate}[label=\alph*)]
\item If \term{CDS spread $<$ bond yield spread}: it is profitable to buy the
bond and hedge with CDS, earning more than the risk-free rate.
\item If \term{CDS spread $>$ bond yield spread}: selling the bond and CDS
protection can be advantageous, effectively borrowing at below the risk-free
rate.
\end{enumerate}

\subsection*{Common reasons for a difference in the CDS-bond basis}

\begin{enumerate}
\item Bonds may sell for much higher than par (negative basis) or lower than par
(positive basis).
\item CDSs are subject to counterparty risk (negative basis).
\item A cheapest-to-deliver possibility exists with a CDS (positive basis).
\item CDS payoffs exclude accrued interest (negative basis).
\item CDS contracts may contain a restructuring clause that allows for payoffs
even without default (positive basis).
\item The bond yield spread is computed using a risk-free rate that is dissimilar
to the one used by the market.
\item \term{Market illiquidity.} Difficulty in exploiting arbitrage opportunities
due to illiquidity can prevent the basis from reaching zero.
\end{enumerate}

\begin{crtrapbox}[The signs are the examinable part]
Five of the seven reasons carry a stated direction, and they alternate. Negative
basis: bonds above par, counterparty risk on the CDS, accrued interest excluded
from the CDS payoff. Positive basis: bonds below par, cheapest-to-deliver
optionality, restructuring clauses. Anything that makes CDS protection
\emph{worth less} than it looks pushes the basis negative; anything that hands
the protection buyer an extra option pushes it positive.
\end{crtrapbox}

% =====================================================================
\lo{CR-9 j}{Compare default probabilities calculated from historical data with those calculated from credit yield spreads.}

\begin{crkeybox}[Three findings]
\begin{itemize}
\item The hazard rates computed from credit spreads are significantly
\textbf{greater} than those computed from historical data.
\item As credit quality improves (worsens), the differences are less (greater).
\item A good (bad) economy decreases (increases) default probabilities.
\end{itemize}
\end{crkeybox}

% =====================================================================
\lo{CR-9 k}{Describe the difference between real-world and risk-neutral default probabilities and determine which one to use in the analysis of credit risk.}

\begin{center}
\footnotesize
\begin{tabular}{L{7.0cm} L{7.5cm}}
\toprule
\thead{Risk-neutral estimates} & \thead{Real-world estimates} \\
\midrule
Derived from financial market prices & Obtained from historical data \\
Estimated from credit spreads & Can rely on credit ratings \\
Used for \textbf{valuation} & Used for \textbf{scenario analysis} \\
\bottomrule
\end{tabular}
\end{center}
\figcap{The decision rule sits in the bottom row. If the question is what
something is worth, use risk-neutral. If the question is what is likely to
happen, use real-world.}

% =====================================================================
\lo{CR-9 l}{Calculate the value of a firm's debt and equity, the volatility of firm value, and the volatility of firm equity using the Merton model.}

The Merton model is a structural approach to predicting default. It assumes
default occurs when the proprietary structure of the defaulting company is no
longer worthwhile. It is based on Black-Scholes-Merton option pricing theory and
evaluates various components of firm value. The terminal payoffs and the
equity-as-a-call picture are set out at \textbf{CR-5 e}; this objective adds the
step that makes the model usable.

\begin{crfmlbox}[Value of equity today]
\[ E_0 \;=\; V_0 N(d_1) - D e^{-rT} N(d_2) \]
\[ d_1 = \frac{\ln(V_0/D) + (r + \sigma_V^{2}/2)T}{\sigma_V\sqrt{T}}
   \qquad d_2 = d_1 - \sigma_V\sqrt{T} \]
\end{crfmlbox}

\gapbadge

{\footnotesize\color{FRMGrey}The objective asks for the volatility of firm value
and the volatility of firm equity, and for the value of the firm's debt. The
source notes carry the objective heading and the equity formula above, but not
the relationship between the two volatilities, not the simultaneous solve, and not
a worked example. The following is written from the GARP curriculum.}

\begin{crgapbox}[What the objective asks for: the two unknowns and how they are found]
Under Merton's model the company defaults when the option is not exercised, and
the probability of this is $N(-d_2)$. To calculate it we require $V_0$ and
$\sigma_V$. \term{Neither of these is directly observable.}

\smallskip
However, if the company is publicly traded we can observe $E_0$, so the equity
equation provides \emph{one} condition that $V_0$ and $\sigma_V$ must satisfy. We
can also estimate $\sigma_E$. From Ito's lemma,
\[ \sigma_E E_0 \;=\; \frac{\partial E}{\partial V}\,\sigma_V V_0
   \;=\; N(d_1)\,\sigma_V V_0 \]
since $\partial E/\partial V$ is the delta of the equity, which is $N(d_1)$.

\smallskip
This provides a \emph{second} condition. The two equations form a pair of
simultaneous equations that can be solved for $V_0$ and $\sigma_V$, in practice
numerically.

\smallskip
The market value of the debt then follows as $V_0 - E_0$.
\end{crgapbox}

\begin{crexbox}[Merton model: solving for firm value and asset volatility]
The value of a company's equity is \$3 million and the volatility of the equity is
80\%. The debt that will have to be paid in one year is \$10 million. The
risk-free rate is 5\% per annum.
\tcblower
Here $E_0 = 3$, $\sigma_E = 0.80$, $r = 0.05$, $T = 1$, and $D = 10$.

\smallskip
Solving the two simultaneous equations yields
\[ V_0 = 12.40 \qquad\text{and}\qquad \sigma_V = 0.2123 \]
The parameter $d_2$ is $1.1408$, so the probability of default is
\[ N(-d_2) = 0.127 \text{, or } 12.7\% \]
The market value of the debt is $V_0 - E_0 = 12.40 - 3 = 9.40$.

The present value of the promised payment on the debt is
$10e^{-0.05\times1} = 9.51$.

The expected loss on the debt is therefore
\[ \frac{9.51 - 9.40}{9.51} \approx 1.2\% \text{ of its no-default value} \]

\smallskip
Note the scale of the difference between the two volatilities: equity volatility
of 80\% corresponds to asset volatility of only 21.23\%, because the equity is a
levered claim on the assets.
\end{crexbox}

% =====================================================================
\lo{CR-9 m}{Calculate distance to default and default probability using the Merton model.}

Distance to default measures the distance between a firm's current financial
position and the point of default. DtD is the number of standard deviations the
company's assets are away from the face value of the debt.

\begin{crfmlbox}[Distance to default and the two probabilities]
\[ \mathrm{DtD} \;=\; d_2 \;=\;
   \frac{\ln(V_0) - \ln(D) + (r - \sigma_V^{2}/2)T}{\sigma_V\sqrt{T}} \]
\[ \text{Risk-neutral PD} = N(-d_2) = 1 - N(d_2) \]
\vk{For the real-world probability of default, the expected return on the firm's
assets $(\mu)$ is substituted for the risk-free rate $(r)$. The risk-neutral
version assumes the firm's assets grow at the risk-free rate.}
\end{crfmlbox}

\begin{crkeybox}[Direction]
The higher (lower) the distance to default, the lower (higher) the probability of
default. $N(-d_2)$ is effectively the probability that shareholders will
\emph{not} exercise their option to repay the loan.
\end{crkeybox}

% =====================================================================
\lo{CR-9 n}{Assess the quality of the default probabilities produced by the Merton model, the Moody's KMV model, and the Kamakura model.}

\begin{enumerate}
\item \term{Merton model.} Uses an option pricing model, and is only applicable
to liquid, publicly traded companies because of the data requirements. Assumes a
standardized normal distribution for default frequencies. Considers only
short-term debt in the default point.
\item \term{Moody's-KMV Expected Default Frequency (EDF) model.} Differs from
Merton in three ways: it uses historical data to develop a distribution of default
frequencies, which is more realistic; it includes short-term debt \emph{and half
of long-term debt}, which better reflects loan obligations; and it can handle
portfolios beyond just debt contracts, such as swaps and convertible debt.
\item \term{CreditMetrics.} Uses credit ratings as an input, which are readily
available for many companies, and considers historical default data through a
transition matrix.
\end{enumerate}

\gapbadge

{\footnotesize\color{FRMGrey}The objective names Merton, Moody's KMV and
\emph{Kamakura}. The source notes cover Merton and Moody's KMV, substitute
CreditMetrics for Kamakura, and do not state the overall assessment the objective
asks for. The following is written from the GARP curriculum.}

\begin{crgapbox}[What the objective asks for: the assessment, and Kamakura]
\textbf{How good are the probabilities?} Merton's model and its extensions
produce a \emph{good ranking} of default probabilities, risk-neutral or
real-world. That is the claim the evidence supports, and it is weaker than a claim
of accuracy in levels.

\smallskip
Because the ranking is good, a \term{monotonic transformation} can be estimated
to convert the PD output from Merton's model into a good estimate of either the
real-world or the risk-neutral default probability.

\smallskip
\textbf{Where Kamakura comes in.} Both \term{Moody's KMV} and \term{Kamakura}
provide a service that transforms a default probability produced by Merton's
model into a \emph{real-world} default probability.

\smallskip
This may seem strange, since $N(-d_2)$ is in theory a \emph{risk-neutral} default
probability, being calculated from an option pricing model. Given the nature of
the calibration process, the underlying assumption is that the rankings of
risk-neutral default probabilities, real-world default probabilities, and the
default probabilities produced by Merton's model are \textbf{all the same}.
\end{crgapbox}
